% Part: computability
% Chapter: recursive-functions
% Section: trees

\documentclass[../../../include/open-logic-section]{subfiles}

\begin{document}

\olfileid{cmp}{rec}{tre}
\olsection{Wit}

Kadhang migunani yen wit diwakili minangka wilangan asli,
kaya rerangken kang bisa diwakili nganggo wilangan, dene
sipat lan operasi ing rerangken mau diwakili dening relasi
lan fungsi rekursif primitif ing kode-kodene. Kita bakal
nggunakake rerangken lan kode-kodene kanggo nindakake iki.
Wit bisa awujud siji simpul (kang bisa nduwé label),
utawa siji simpul (kang bisa nduwé label) kang
disambungake karo sawetara subwit. Simpul mau diarani
\emph{oyod} wit, lan subwit kang langsung kasambung
diarani \emph{subwit langsung}.

Kita ngode wit kanthi rekursif minangka rerangken
$\tuple{k, d_1, \dots, d_k}$, kanthi $k$ cacah subwit
langsung lan $d_1$, \dots,~$d_k$ kode subwit
langsung kasebut. Yen simpul-simpule mawa label,
label kasebut bisa dilebokake sawisé subwit langsung.
Dadi wit kang mung ngemot siji simpul mawa label~$l$
dikode nganggo $\tuple{0,l}$, dene wit kang ngemot
oyod (mawa label~$l_1$) kang kasambung karo rong
simpul tunggal (mawa label $l_2$, $l_3$)
dikode nganggo $\tuple{2,
  \tuple{0, l_2}, \tuple{0, l_3}, l_1}$.

\begin{prop}
  \ollabel{prop:subtreeseq}
  Fungsi $\fn{SubtreeSeq}(t)$, kang ngasilake kode
  rerangken kanthi unsur-unsur awujud kode kabeh
  subwit saka wit kang kode~$t$, iku rekursif primitif.
\end{prop}

\begin{proof}
  Pisanan, gatekna yen $\fn{ISubtrees}(t) =
  \fn{subseq}(t, 1, (t)_0)$ iku rekursif primitif
  lan ngasilake kode subwit langsung saka wit~$t$.
  Saiki kita bisa netepake fungsi pambantu
  $\fn{hSubtreeSeq}(t,n)$ kang ngitung rerangken
  kabeh subwit kang jarake saka oyod paling akeh
  $n$ rusuk, kanthi kemungkinan ana ulangan.
  Rerangken subwit saka~$t$ kang jarake $0$
  rusuk saka oyod---yaiku diwiwiti ing oyod
  wit~$t$---mung ngemot~$t$.
  Kanggo nggayuh tataran nganti~$n+1$ saka $t$, kita
  nyambungake rerangken kang wis ngemot subwit
  nganti tataran $n$ karo rerangken kabeh subwit
  langsung saka subwit nganti tataran $n$.
  % OLPL-131: Rekurensi n+1 nambah subwit marang kabeh asil nganti n;
  % asil iki kumulatif, dudu mung subwit persis tataran n.
  Kanggo ndhaptar kabeh subwit langsung iki, gatekna
  yen $f(x)$ iku fungsi rekursif primitif kang
  ngasilake kode rerangken, mula $g_f(s, k) = f((s)_0) \concat
  \dots \concat f((s)_k)$ uga rekursif primitif:
    \begin{align*}
      g(s, 0) & = f((s)_0)\\
      g(s, k+1) & = g(s, k) \concat f((s)_{k+1})
    \end{align*}
    Umpamane, yen $s$ rerangken wit, mula
    $h(s) = g_{\fn{ISubtrees}}(s, \len{s})$
    ngasilake rerangken subwit langsung saka
    unsur-unsur~$s$. Suku tambahan kang indeksé
    padha karo dawane rerangken mung ngasilake
    rerangken kosong: unsur ing njaban dawa
    diwenehi nilai nol, lan subwit langsung saka
    kode kosong uga kosong.
    % OLPL-132: g(s,k) ngemot indeks 0 nganti k;
    % nalika k=dawa s, suku tambahan f(0) iku kosong.
    Kita bisa nggunakake fungsi iki kanggo netepake
    $\fn{hSubtreeSeq}$ kanthi
    \begin{align*}
      \fn{hSubtreeSeq}(t, 0) & = \tuple{t} \\
      \fn{hSubtreeSeq}(t, n+1) & = \fn{hSubtreeSeq}(t, n) \concat
      h(\fn{hSubtreeSeq}(t, n)).
    \end{align*}
    Tataran paling gedhe saka subwit ing wit kang
    dikode nganggo~$t$, yaiku jarak paling gedhe
    antarane oyod lan simpul godhong, diwatesi
    dening kode~$t$. Mula rerangken kode kabeh
    subwit saka wit kang dikode nganggo~$t$
    diwenehake dening $\fn{hSubtreeSeq}(t, t)$.
\end{proof}

\begin{prob}
  Definisi $\fn{hSubtreeSeq}$ ing bukti
  \olref[cmp][rec][tre]{prop:subtreeseq} lumrahe
  ngemot ulangan. Wenehana definisi liya kang
  njamin yen kode saben subwit mung katon sepisan
  ing dhaptar asile.
\end{prob}

\end{document}
